\documentclass[11pt,a4paper]{article}

% ============================================================
%  S-JEPA : explication complete et detaillee + FAQ
%  Document pedagogique accompagnant la reimplementation
%  de production (Python = entrainement, Rust = inference).
% ============================================================

\usepackage[utf8]{inputenc}
\usepackage[T1]{fontenc}
\usepackage[french]{babel}
\usepackage{lmodern}
\usepackage{microtype}
\usepackage[margin=2.3cm]{geometry}

\usepackage{amsmath}
\usepackage{amssymb}
\usepackage{amsfonts}
\usepackage{mathtools}

\usepackage{graphicx}
\usepackage{booktabs}
\usepackage{multirow}
\usepackage{array}
\usepackage{longtable}
\usepackage{xcolor}
\usepackage{tcolorbox}
\usepackage{listings}
\usepackage{enumitem}

\usepackage[colorlinks=true,linkcolor=blue!60!black,urlcolor=blue!50!black,
            citecolor=blue!50!black]{hyperref}

% --- Couleurs / boites ---
\definecolor{warnbg}{RGB}{255,247,230}
\definecolor{warnframe}{RGB}{200,120,0}
\definecolor{faqbg}{RGB}{238,245,255}
\definecolor{faqframe}{RGB}{40,90,170}
\definecolor{codebg}{RGB}{245,245,245}

\newtcolorbox{warnbox}{colback=warnbg,colframe=warnframe,boxrule=0.6pt,
  arc=2pt,left=6pt,right=6pt,top=4pt,bottom=4pt}

\newtcolorbox{faqbox}{colback=faqbg,colframe=faqframe,boxrule=0.6pt,
  arc=2pt,left=6pt,right=6pt,top=4pt,bottom=4pt}

% --- Style code ---
\lstset{
  basicstyle=\ttfamily\footnotesize,
  backgroundcolor=\color{codebg},
  breaklines=true,
  frame=single,
  rulecolor=\color{gray!40},
  showstringspaces=false,
  columns=fullflexible,
  keepspaces=true,
}

% --- Macro FAQ ---
\newcommand{\faq}[2]{%
  \par\medskip
  \begin{faqbox}
  \textbf{Q. #1}\par\smallskip
  \textbf{R.} #2
  \end{faqbox}
}

\newcommand{\code}[1]{\texttt{#1}}
\newcommand{\file}[1]{\texttt{#1}}

\title{\textbf{S-JEPA : explication complète et détaillée}\\[0.4em]
       \large \emph{Soft Clustering Anchors for Self-Supervised Speech
       Representation Learning}\\[0.6em]
       \normalsize Document pédagogique pour la réimplémentation de production}
\author{Notes de travail --- dépôt \texttt{sjepa}}
\date{\today}

\begin{document}
\maketitle

\begin{abstract}
\noindent
Ce document explique de bout en bout le papier \textbf{S-JEPA}
(arXiv:2606.19398) et le codebase de référence fourni
(\file{model.py}, \file{train.py}, \file{fit\_gmm.py}). Il relie
chaque concept mathématique à sa ligne de code, recense les écarts
entre le papier et le code de référence (importants pour une
réimplémentation fidèle), puis répond à une large série de questions
« bêtes mais légitimes » dans une FAQ. Objectif du dépôt :
réimplémenter S-JEPA pour la production --- \textbf{entraînement en
Python}, \textbf{inférence en Rust}.
\end{abstract}

\tableofcontents
\bigskip\hrule

% ============================================================
\section{Vue d'ensemble en une page}
% ============================================================

\textbf{Le problème.} L'apprentissage auto-supervisé (SSL) de la parole
est dominé par une recette unique (HuBERT, WavLM) : on clusterise des
features acoustiques \emph{hors-ligne} par $k$-means, chaque frame reçoit
un \emph{identifiant de cluster unique} (cible « dure »), puis on entraîne
un encodeur à prédire cet ID aux positions masquées via une
cross-entropy. On interrompt ensuite l'entraînement pour re-clusteriser
tout le corpus sur des features intermédiaires, et on recommence.

\textbf{Les deux défauts attaqués.}
\begin{enumerate}[leftmargin=*]
  \item \textbf{Assignation dure} : aux frontières de phonèmes, transitions
  et silences, l'ambiguïté acoustique réelle est écrasée dans une seule
  catégorie.
  \item \textbf{Pipeline « stop-restart »} : re-clusteriser le corpus entier
  entre itérations est coûteux et interrompt l'entraînement.
\end{enumerate}

\textbf{L'idée de S-JEPA.}
\begin{itemize}[leftmargin=*]
  \item \textbf{Cibles douces} : la cible n'est plus un ID unique mais la
  \emph{distribution postérieure complète d'un GMM} (Gaussian Mixture Model)
  sur $K$ composantes.
  \item \textbf{Loss = KL divergence} entre la prédiction et le postérieur GMM.
  \item \textbf{GMM en ligne} en phase 2 : le GMM se met à jour en continu,
  ce qui supprime le re-clustering hors-ligne.
  \item \textbf{Architecture JEPA} : encodeur + prédicteur + tête de cluster,
  avec un encodeur EMA pour alimenter le GMM en phase 2.
\end{itemize}

\textbf{Le résultat.} Avec un encodeur de 51.8M de paramètres, S-JEPA est
la méthode SSL au plus faible WER sous 90M paramètres
(\textbf{12.10\,\%} sur LibriSpeech \texttt{test-clean}, greedy CTC ;
8.50\,\% avec LM) et égale HuBERT-Base en reconnaissance d'émotion à 55\,\%
de ses paramètres.

\begin{warnbox}
\textbf{À retenir tout de suite.} Le codebase fourni ne contient que le
\emph{pré-entraînement}. Il n'y a ni évaluation SUPERB, ni extraction de
features, ni inférence. C'est exactement ce que notre projet doit
construire (cf. §\ref{sec:gaps} et §\ref{sec:prod}).
\end{warnbox}

% ============================================================
\section{Intuition : cibles dures contre cibles douces}
% ============================================================

Imaginons une frame audio située pile à la frontière entre deux sons,
par exemple entre le « s » et le « t » de « \emph{st}op ». Acoustiquement,
cette frame est \emph{ambiguë} : elle ressemble à 45\,\% au cluster « s »,
à 40\,\% au cluster « t », et un peu à d'autres.

\begin{itemize}[leftmargin=*]
  \item \textbf{Cible dure (HuBERT)} : on force la frame dans \emph{une}
  classe, disons « s ». L'information « cette frame est aussi presque un t »
  est \emph{définitivement perdue}.
  \item \textbf{Cible douce (S-JEPA)} : la cible est le vecteur
  $[0.45, 0.40, 0.05, \dots]$. Le modèle apprend à reproduire toute la
  distribution, donc \emph{l'ambiguïté est préservée}.
\end{itemize}

Le papier mesure ceci directement (§\ref{sec:analysis}) : sur les frames
de parole réelles, l'entropie de la distribution prédite est
\emph{bimodale} --- soit très confiante, soit en « égalité à deux » ---,
et plus d'un tiers des frames portent plus d'un bit d'incertitude. C'est
la preuve empirique que les cibles douces capturent une information que
les cibles dures jetteraient.

% ============================================================
\section{La loss : KL entre prédiction et postérieur GMM}
% ============================================================

\subsection{Le postérieur d'un GMM diagonal}

Un GMM à $K$ composantes et covariance \emph{diagonale} donne, pour un
vecteur de features $m$, la probabilité a posteriori d'appartenir à la
composante $k$ :
\begin{equation}
q_k(m) = \frac{\pi_k\,\mathcal{N}(m;\mu_k,\sigma_k^2)}
              {\sum_j \pi_j\,\mathcal{N}(m;\mu_j,\sigma_j^2)}.
\end{equation}
En pratique on le calcule en espace log pour la stabilité numérique :
\begin{equation}
q_k(m) = \mathrm{softmax}_k\!\left(
  \log\pi_k - \frac{D}{2}\log(2\pi)
  - \frac{1}{2}\sum_{d=1}^{D}\Big[
      \log\sigma_{k,d}^2 + \frac{(m_d-\mu_{k,d})^2}{\sigma_{k,d}^2}
    \Big]\right).
\label{eq:gmm_post}
\end{equation}
C'est \emph{exactement} ce qu'implémentent \code{GMMAssigner.posteriors}
(\file{train.py:63}) et \code{OnlineGMM.posteriors} (\file{train.py:111}) :
le terme $\sum_d (m_d-\mu_{k,d})^2/\sigma_{k,d}^2$ est la distance de
Mahalanobis (\code{mahal}), $\sum_d\log\sigma_{k,d}^2$ est \code{log\_det},
puis \code{logsumexp} normalise. Les variances sont « clampées » à
$10^{-6}$. Le calcul est découpé en blocs (\emph{chunks}) de 500 pour
borner la mémoire.

\subsection{La fonction de coût}

La loss est la divergence de Kullback--Leibler entre le postérieur GMM
$q_t$ (la \emph{cible}, calculée sous \code{no\_grad}) et la sortie
softmax du prédicteur $p_t$, moyennée sur un ensemble de positions $S$ :
\begin{equation}
\mathcal{L} = \frac{1}{|S|}\sum_{t\in S}\mathrm{KL}\big(q_t \,\|\, p_t\big),
\qquad p_t = \mathrm{softmax}\big(g_\omega(\hat z_t)\big).
\label{eq:loss}
\end{equation}
Dans le code (\file{train.py:833--841}) : \code{log\_pred\_mask =
log\_softmax(logits)}, puis \code{F.kl\_div(log\_pred\_mask, target\_mask,
reduction='batchmean')}. Seules les positions masquées contribuent au
gradient (\code{loss = loss\_mask}) ; la KL sur positions visibles est
calculée \emph{sous \code{no\_grad}} uniquement pour le logging.

\begin{warnbox}
\textbf{Sens de la KL et ordre des arguments.}
\code{torch.nn.functional.kl\_div(input, target)} attend
\code{input} en \emph{log}-probabilités et calcule
$\sum target\cdot(\log target - input)$, soit $\mathrm{KL}(target\,\|\,
\exp(input)) = \mathrm{KL}(q\,\|\,p)$. C'est bien la « forward KL » du
papier : la cible $q$ (GMM) pondère l'erreur. À reproduire à l'identique
en Rust pour garantir la parité numérique.
\end{warnbox}

% ============================================================
\section{Architecture détaillée (\texorpdfstring{$\leftrightarrow$}{<->} \file{model.py})}
% ============================================================

L'architecture suit le pattern JEPA avec \textbf{trois composants
entraînés} --- encodeur $f_\phi$, prédicteur $h_\psi$, tête de cluster
$g_\omega$ --- plus un \textbf{encodeur EMA} $\bar f$ non entraîné
(utilisé seulement en phase 2).

\subsection{Le flux de données}

\begin{lstlisting}
waveform [B,1,T]
  v
ConvFeatureExtractor   7 blocs conv, stride total 320 -> frames 20 ms (50 Hz)
  |  (model.py:30  kernels (10,5)(3,2)x4(2,2)x2, GELU, GroupNorm fp32 au 1er bloc)
post_extract_proj (512->768) -> Fp32LayerNorm
  |
zerote les frames masquees  (l'encodeur ne voit RIEN aux positions masquees)
  |  (model.py:209  x = x * mask.unsqueeze(-1))
ConvPositionalEncoding (kernel 128, 16 groupes, weight-norm) -> encoder LayerNorm
  |
N x TransformerLayer (MHSA + FFN, post-norm, LayerNorm fp32)
  |  z = f_phi(x)
  +----------------------------+
  |                            |
cluster_head(z)          JEPAPredictor : injecte mask_token, ajoute pos_embed,
-> logits "visibles"     1 couche Transformer -> z_hat
(logges, NON supervises) puis cluster_head(z_hat) -> logits "masques"
                              |
                     KL( logits masques || posteriors GMM )   <- SEUL signal
\end{lstlisting}

\subsection{Les composants un par un}

\begin{itemize}[leftmargin=*]
  \item \textbf{\code{ConvFeatureExtractor}} (\file{model.py:30}) : 7 conv
  1D, kernels/strides identiques à HuBERT, stride total $=320$ échantillons
  $\Rightarrow$ frames de 20\,ms à 16\,kHz (50\,Hz). GroupNorm fp32 au
  premier bloc, GELU partout.

  \item \textbf{\code{GradMultiply}} (\file{model.py:68}) : multiplie le
  gradient du frontend CNN par \code{feature\_grad\_mult=0.1} pour
  empêcher le frontend d'apprendre trop vite (repris de HuBERT).

  \item \textbf{\code{ConvPositionalEncoding}} (\file{model.py:54}) :
  encodage positionnel \emph{convolutionnel} (kernel 128, 16 groupes,
  \code{weight\_norm}) ajouté avant la pile de transformers.

  \item \textbf{\code{TransformerLayer}} (\file{model.py:78}) : attention
  multi-têtes + FFN, \textbf{post-norm}, LayerNorm forcé en fp32
  (\code{Fp32LayerNorm}) pour la stabilité en précision mixte.

  \item \textbf{\code{OnlineEncoder}} (\file{model.py:169}) : assemble tout.
  Contient aussi \code{cluster\_head} (MLP $768\!\to\!768\!\to\!768\!\to\!K$)
  et \code{predictor}.

  \item \textbf{\code{JEPAPredictor}} (\file{model.py:119}) : \code{mask\_token}
  apprenable injecté aux positions masquées, \code{pos\_embed} (capacité 750
  frames $=15$\,s), \textbf{une seule} couche \code{TransformerEncoderLayer}
  (8 têtes, FFN $=1536$), puis \code{proj}. Il ne remplit \emph{que} les
  positions masquées.

  \item \textbf{\code{cluster\_head}} : MLP partagé entre chemin visible et
  masqué, mais seul le chemin masqué entre dans la loss.
\end{itemize}

\subsection{Convention de masque : attention au piège}

\begin{warnbox}
Dans le \emph{code}, \code{mask} vaut \textbf{1 = visible, 0 = masqué}
(\file{train.py:812} : \code{mask = (\textasciitilde mask\_bool).float()}),
ce qui est l'\emph{inverse} de la notation du papier. L'encodeur fait
\code{x = x * mask.unsqueeze(-1)} : il met donc à zéro les positions
\emph{masquées}. Le \code{mask\_token} n'est injecté que plus tard, par le
prédicteur. À fixer comme convention unique et documentée dans notre
réimplémentation.
\end{warnbox}

% ============================================================
\section{Les cibles : le GMM en phase 1 et en phase 2}
% ============================================================

\subsection{Phase 1 --- GMM figé sur MFCC (\file{fit\_gmm.py})}

\begin{itemize}[leftmargin=*]
  \item \textbf{Features} : 39-dim MFCC = 13 coefficients cepstraux $+\ \Delta\ +\
  \Delta\Delta$ (\code{extract\_mfcc}, \file{fit\_gmm.py:254}), frames 20\,ms.
  \item \textbf{Reservoir sampling} (\code{collect\_features},
  \file{fit\_gmm.py:266}) : échantillonne les frames en flux sans tout charger
  en mémoire (algorithme de Vitter).
  \item \textbf{Init $k$-means} (\code{gmm\_init\_kmeans},
  \file{fit\_gmm.py:313}, 5 itérations) puis \textbf{EM}
  (\code{fit\_gmm\_em}, \file{fit\_gmm.py:378}, 20 itérations). Statistiques
  suffisantes accumulées par chunks ; jamais la matrice $N\times K$ entière.
  \item \textbf{Sortie} : \file{gmm.pt} $=\{$means, variances, weights, $K$,
  dim$\}$. $K=100$ dans le papier.
  \item Le GMM est ensuite \textbf{figé} pendant toute la phase 1.
\end{itemize}

\subsection{Phase 2 --- GMM en ligne sur features de l'encodeur EMA}

\code{OnlineGMM} (\file{train.py:90}) :
\begin{itemize}[leftmargin=*]
  \item \textbf{Features} : 768-dim, sortie de l'encodeur EMA $\bar f$ à la
  couche active $\ell^\star$.
  \item \textbf{Mise à jour} (\code{OnlineGMM.update}, \file{train.py:134}) :
  à chaque batch on calcule des moyennes/variances/poids \emph{pondérés par
  les responsabilités} $q_{n,k}$, puis une mise à jour EMA :
  \begin{equation}
  \theta \leftarrow \alpha\,\theta + (1-\alpha)\,\hat\theta,
  \qquad \theta \in \{\pi,\mu,\sigma^2\},\quad \alpha = \texttt{param\_ema} = 0.999.
  \end{equation}
  \item \textbf{Aucun gradient ne traverse le GMM.} Coût $O(B\cdot K\cdot D)$
  par pas, intégré dans la passe forward $\Rightarrow$ pas de re-clustering
  hors-ligne.
  \item \textbf{Seeding} (\file{train.py:610--737}) : au début de la phase 2,
  le GMM est initialisé par un mini-batch $k$-means (Sculley 2010) sur les
  features de l'encodeur EMA, puis broadcasté à tous les rangs.
\end{itemize}

% ============================================================
\section{L'entraînement en deux phases}
% ============================================================

Les deux phases forment \textbf{une seule trajectoire d'optimisation
continue}, mais sont lancées comme \textbf{deux exécutions} du même
\file{train.py}.

\begin{center}
\begin{tabular}{@{}p{0.46\textwidth}p{0.46\textwidth}@{}}
\toprule
\textbf{Phase 1 (GMM MFCC figé)} & \textbf{Phase 2 (GMM en ligne)} \\
\midrule
\code{--gmm\_path gmm\_mfcc.pt}, sans \code{--online\_gmm} &
\code{--online\_gmm --reseed\_gmm --K 500}, repart de \code{--pt\_ckpt} \\
Cibles = GMM MFCC figé &
Cibles = GMM en ligne sur features EMA \\
Loss sur masqué \emph{+} visible &
Loss masqué+visible (début) puis masqué seul \\
Augmentation activée ($p=0.25$) &
Augmentation activée puis désactivée \\
$K=100$, LR $=10^{-4}$ &
$K=500$, LR $=2.5\times10^{-5}$ \\
\bottomrule
\end{tabular}
\end{center}

% ============================================================
\section{Les trois innovations de la phase 2}
% ============================================================

Ce sont les vraies contributions techniques : elles remplacent le
stop-restart de HuBERT par un entraînement continu.

\subsection{(a) GMM en ligne}
Décrit ci-dessus. Supprime la passe de re-labellisation du corpus entier
entre itérations.

\subsection{(b) Sélection adaptative de couche par \emph{effective rank}}

Quelle couche du transformer alimente le GMM ? HuBERT la choisit à la
main ; S-JEPA l'automatise via le \textbf{rang effectif} :
\begin{equation}
\mathrm{erank}(Z) = \exp\!\Big(-\sum_i p_i\log p_i\Big),
\qquad p_i = \frac{\sigma_i}{\sum_j \sigma_j},
\end{equation}
où les $\sigma_i$ sont les valeurs singulières d'un batch de 2000
embeddings centrés. C'est l'exponentielle de l'entropie de Shannon du
spectre : élevé quand beaucoup de valeurs singulières sont comparables
(représentation « riche »). Implémentation : \file{train.py:886--912}
(\code{torch.linalg.svdvals}, normalisation, entropie, lissage EMA, puis
\code{argmax}). Observé : $\ell^\star=2$ en fin de phase 1, puis transition
vers $\ell^\star=4$ au milieu de la phase 2.

\subsection{(c) Décroissance EMA périodiquement commutée (\emph{switched EMA})}

Tension fondamentale : l'encodeur EMA \emph{alimente} le GMM, dont les
composantes sont estimées sur ces mêmes features.
\begin{itemize}[leftmargin=*]
  \item EMA trop lente $\Rightarrow$ le GMM ne voit jamais les améliorations
  récentes de l'encodeur.
  \item EMA trop rapide $\Rightarrow$ la distribution cible bouge à chaque pas.
\end{itemize}
Solution : alterner entre deux valeurs fixes,
\begin{equation}
\alpha_t \in \{\alpha_{\text{fast}}=0.999,\ \alpha_{\text{slow}}=0.9999\},
\qquad \bar\phi \leftarrow \alpha_t\,\bar\phi + (1-\alpha_t)\,\phi,
\end{equation}
avec une cadence moyenne d'environ 20\,000 pas. C'est l'\textbf{analogue
continu} du re-cluster hors-ligne : phase rapide = « le GMM rattrape
l'encodeur » ; phase lente = « cible stable, l'encodeur apprend contre
elle ». Échelles caractéristiques : $\tau(\alpha)=1/(1-\alpha)$, soit
$1\,000$ pas (fast) et $10\,000$ pas (slow).

% ============================================================
\section{L'algorithme complet (Algorithm 1)}
% ============================================================

Correspondance directe avec la boucle \file{train.py:746--960} :

\begin{lstlisting}
Init: f_phi, h_psi, g_omega, encodeur EMA f_bar, GMM, couche ell*, alpha <- alpha_fast
Pour chaque minibatch x :
   x_aug, x_clean <- Augment(x)                          # train.py:781
   m <- MFCC(x_clean)  [Ph.1]  ou  f_bar(x_clean) a ell* [Ph.2]   (no_grad)
   q <- GMM.posterior(m)                                 # :795 / :806
   si Phase 2 :
       GMM <- EMA-update depuis m, q                     # :860 (online, no grad)
       si pas de verif-couche : ell* <- argmax_l erank(Z^l)  # :886-912
   mask <- SampleBlockMask()                             # :811
   z <- f_phi(x_aug, mask) ;  z_hat <- h_psi(z, mask)    # :815
   p_t <- softmax(g_omega(z_hat_t)) pour t dans M
   L <- (1/|M|) somme KL(q_t || p_t) ; backward ; step   # :837-849
   si Phase 2 : flip alpha ; f_bar <- alpha f_bar + (1-alpha) f_phi  # :856-857
Retour: f_phi   (predicteur, cluster_head, GMM jetes)
\end{lstlisting}

\textbf{Masquage par blocs} (\code{compute\_mask\_indices},
\file{train.py:362}) : on tire \code{num\_spans} spans de 10 frames jusqu'à
viser un ratio de 0.65. La variante rapide ne garantit pas exactement
0.65 (chevauchements possibles) ; une variante exacte \code{\_slow} existe
en \file{train.py:348}.

% ============================================================
\section{Infrastructure d'entraînement (orientée production)}
% ============================================================

\begin{itemize}[leftmargin=*]
  \item \textbf{DeepSpeed} pour le multi-GPU (8 GPU, batch global 192),
  config \file{ds\_config\_jepa\_pretrain.json}.
  \item \textbf{Données} : manifests JSONL \code{\{"wav\_path": ...\}}, lus en
  streaming (\code{StreamingDataset}, \file{train.py:243}). Un
  \file{path\_index.txt} globalement mélangé est construit une seule fois sur
  rank 0 ; chaque worker lit des slices disjointes (\code{i \% total == wid}).
  \item \textbf{Checkpoint portable} (\file{portable\_step*.pt},
  \file{train.py:917--942}) : bundle \{encodeur online, encodeur EMA, params
  GMM, step\}, agnostique du framework, écrit dans un thread daemon pour ne
  pas bloquer l'entraînement.
  \item \textbf{Augmentation} (\code{DenoiseAugmentor}, \file{train.py:190}) :
  noise mixing au SNR + utterance mixing, suivant WavLM. Le GMM voit le
  signal \emph{propre}, l'encodeur voit le signal \emph{augmenté}.
\end{itemize}

% ============================================================
\section{Évaluation et résultats}
% ============================================================

Protocole \textbf{SUPERB} : encodeur \emph{gelé}, on entraîne de petites
têtes par tâche.

\begin{center}
\begin{tabular}{@{}lll@{}}
\toprule
\textbf{Tâche} & \textbf{S-JEPA (51.8M)} & \textbf{Comparaison} \\
\midrule
ASR (WER test-clean) & \textbf{12.10\,\%} greedy ; 8.50\,\% +LM
  & meilleur SSL $<$90M (DeCoAR 2.0 : 13.02\,\% à 89.8M) \\
Émotion (ER, acc.)   & \textbf{64.83\,\%}
  & $\approx$ HuBERT-Base (64.92\,\%) à 55\,\% des params \\
Slot Filling (F1/CER) & 83.05 / 33.17
  & $\approx$ DeCoAR 2.0 \\
\bottomrule
\end{tabular}
\end{center}

Plus du \textbf{probing} (sondes linéaires/MLP sur représentations gelées) :
S-JEPA domine Speaker ID, Gender, Chapter ID, et gagne en Phoneme avec une
sonde MLP. Conclusion : nouvelle \textbf{frontière de Pareto sous 90M
paramètres}, sans \emph{teacher} ni re-clustering hors-ligne.

% ============================================================
\section{L'analyse centrale : l'entropie bimodale}
\label{sec:analysis}
% ============================================================

Le résultat empirique qui justifie les cibles douces : sur $\sim$153\,000
frames held-out, l'entropie par frame de la prédiction
\begin{equation}
H_t = -\sum_{k=1}^{K} p_{t,k}\log_2 p_{t,k} \in [0, \log_2 K]
\end{equation}
est \textbf{bimodale} :
\begin{itemize}[leftmargin=*]
  \item un mode confiant $< 0.3$ bit (régions phonétiquement stables) ;
  \item un mode secondaire $\approx 1$ bit = égalité parfaite à deux clusters
  (frontières acoustiques) ;
  \item \textbf{36.6\,\%} des frames portent plus de 1 bit.
\end{itemize}
Lecture en bits : $H=1$ bit = égalité 50/50 ; $H=2$ bits = uniforme à
quatre ; $H=\log_2 500\approx 9$ bits = uniforme sur tous les clusters
(jamais atteint). La bimodalité est l'observation « porteuse » : ce n'est
ni du flou diffus (qui serait unimodal), ni un repli dégénéré vers
l'uniforme (qui s'empilerait à $\log_2 K$).

% ============================================================
\section{Écarts entre le papier et le code de référence}
\label{sec:gaps}
% ============================================================

Cruciaux pour une réimplémentation \emph{fidèle au papier}.

\begin{center}
\begin{tabular}{@{}p{0.30\textwidth}p{0.30\textwidth}p{0.30\textwidth}@{}}
\toprule
\textbf{Élément} & \textbf{Papier} & \textbf{Code fourni} \\
\midrule
Couches encodeur & 6 (51.8M) & default 12 ; \code{--num\_layers 6} possible \\
Switched EMA ($\alpha_\text{fast}/\alpha_\text{slow}$)
  & spécifié en détail & \textbf{absent} (\code{--ema\_decay} fixe) \\
Transition loss vis+masq $\to$ masqué seul
  & décrite (annexe) & code n'a \textbf{que} masqué (\code{loss = loss\_mask}) \\
Eval SUPERB / inférence & résultats reportés & \textbf{non fourni} \\
$K$ phase 1 / phase 2 & 100 / 500 & configurable (\code{--K}) \\
\bottomrule
\end{tabular}
\end{center}

% ============================================================
\section{Plan de réimplémentation de production}
\label{sec:prod}
% ============================================================

\textbf{Décision arrêtée :} \emph{entraînement en Python, inférence en
Rust}.

\begin{itemize}[leftmargin=*]
  \item \textbf{Python (entraînement + recherche)} : package modulaire
  reprenant \file{fit\_gmm.py}, \file{train.py}, \file{model.py},
  enrichi des éléments manquants (6 couches, switched EMA, transition de
  loss), testé, avec configs et reprise propre. Production des checkpoints
  portables.
  \item \textbf{Rust (inférence + serving)} : moteur d'inférence
  haute-performance chargeant le checkpoint portable (encodeur seul --- le
  prédicteur, la tête de cluster et le GMM sont jetés après
  pré-entraînement). Pistes : \code{tch} (bindings libtorch) ou
  \code{candle}. Tests de \textbf{parité numérique} contre la sortie
  Python.
  \item \textbf{Évaluation} : à construire (extraction de features gelées +
  têtes SUPERB / probing). Peut rester en Python dans un premier temps.
\end{itemize}

% ============================================================
\section{FAQ --- questions « bêtes » mais légitimes}
% ============================================================

\subsection{Concepts généraux}

\faq{C'est quoi « auto-supervisé » (SSL) au juste ?}{%
On entraîne le modèle \emph{sans étiquettes humaines}. Le signal
d'apprentissage est fabriqué à partir des données elles-mêmes : ici, on
masque des morceaux du signal et on demande au modèle de prédire ce qui
manque (dans un espace de features, pas le son brut). Aucune transcription
n'est nécessaire pour le pré-entraînement.}

\faq{Que signifie « JEPA » ?}{%
\emph{Joint-Embedding Predictive Architecture}. On ne prédit pas les pixels
ou les échantillons audio bruts, mais une \emph{représentation} (embedding)
de la partie masquée. Un encodeur encode le contexte visible, un prédicteur
« remplit » les positions masquées dans l'espace latent.}

\faq{Pourquoi prédire dans l'espace latent plutôt que reconstruire le son ?}{%
Reconstruire le signal brut force le modèle à modéliser des détails non
pertinents (bruit de fond, phase exacte). Prédire une représentation laisse
le modèle se concentrer sur la structure utile (contenu phonétique,
prosodie).}

\faq{C'est quoi un GMM, expliqué simplement ?}{%
Un \emph{Gaussian Mixture Model} suppose que les données sont générées par
un mélange de $K$ « bulles » gaussiennes. Chaque bulle a un centre
($\mu_k$), une largeur ($\sigma_k^2$) et un poids ($\pi_k$). Le
\emph{postérieur} donne, pour un point, la probabilité d'appartenir à
chaque bulle. C'est un $k$-means « mou » qui rend des probabilités au lieu
d'un seul gagnant.}

\faq{Quelle différence concrète entre $k$-means et GMM ici ?}{%
$k$-means assigne chaque point à \emph{un} cluster (cible dure). Le GMM
rend une \emph{distribution} sur tous les clusters (cible douce). C'est
exactement le cœur de S-JEPA : remplacer la cible dure de HuBERT par la
cible douce du GMM.}

\faq{Pourquoi « covariance diagonale » et pas complète ?}{%
Une covariance complète a $D^2$ paramètres par composante (ici $768^2$),
trop coûteux et instable. Diagonale = on suppose les dimensions
indépendantes, $D$ paramètres par composante. Standard pour la parole, et
le postérieur a une forme close simple (Eq.~\ref{eq:gmm_post}).}

\faq{C'est quoi la divergence KL, en une phrase ?}{%
Une mesure d'écart entre deux distributions de probabilité.
$\mathrm{KL}(q\,\|\,p)=0$ si et seulement si $q=p$, et grandit quand $p$
s'éloigne de la cible $q$. Minimiser la KL, c'est rendre la prédiction $p$
aussi proche que possible de la cible $q$.}

\subsection{Architecture et masquage}

\faq{Pourquoi des frames de 20 ms ?}{%
Le frontend CNN a un stride total de 320 échantillons. À 16\,kHz,
$320/16000 = 20$\,ms par frame, soit un débit de 50 frames/seconde. C'est
le standard HuBERT/wav2vec et un bon compromis pour la parole.}

\faq{Pourquoi mettre les frames masquées à zéro \emph{avant} le transformer,
mais injecter le mask token \emph{après} (dans le prédicteur) ?}{%
L'encodeur doit produire une représentation du \emph{contexte visible}
seulement. Si on lui donnait déjà un mask token, il « tricherait ». Le
prédicteur, lui, reçoit le contexte encodé et doit \emph{deviner} le
contenu masqué : c'est là qu'on place le mask token apprenable.}

\faq{Pourquoi multiplier le gradient du CNN par 0.1 ?}{%
Le frontend convolutionnel apprend très vite et peut déstabiliser
l'entraînement. Réduire son gradient ($\gamma_\text{feat}=0.1$) le fait
évoluer plus lentement que les transformers. Astuce héritée de HuBERT.}

\faq{Pourquoi un encodage positionnel \emph{convolutionnel} plutôt que
sinusoïdal ?}{%
Pour la parole, une convolution capture les positions \emph{relatives} de
façon souple et performe mieux empiriquement que les positions absolues
sinusoïdales. C'est le choix de wav2vec~2.0 / HuBERT.}

\faq{Le prédicteur n'a qu'\emph{une} couche transformer, ce n'est pas
trop peu ?}{%
C'est volontaire : le prédicteur est un module \emph{léger et jetable}.
Tout le travail représentationnel doit être fait par l'encodeur, qui est
le seul composant conservé après pré-entraînement.}

\faq{Pourquoi 65\,\% de masquage, c'est énorme, non ?}{%
Oui, et c'est voulu : un masquage élevé rend la tâche difficile et force
l'encodeur à apprendre une vraie structure plutôt qu'une interpolation
locale. 65\,\% avec des spans de 10 frames est la valeur HuBERT.}

\subsection{Encodeur EMA et phase 2}

\faq{C'est quoi l'encodeur EMA et pourquoi en a-t-on besoin ?}{%
C'est une copie « lente » de l'encodeur, dont les poids sont une moyenne
mobile exponentielle de l'encodeur entraîné. En phase 2, il fournit des
features \emph{stables} sur lesquelles le GMM est ajusté. Sans cela, le GMM
poursuivrait une cible qui bouge à chaque pas (instable).}

\faq{Si le GMM dépend de l'encodeur et que l'encodeur apprend contre le
GMM, n'est-ce pas circulaire / risque d'effondrement ?}{%
C'est exactement la tension que gère la \emph{switched EMA} : en alternant
une phase « rapide » (le GMM rattrape l'encodeur) et une phase « lente »
(cible figée, l'encodeur apprend contre elle), on évite que les deux
dérivent ensemble vers une solution dégénérée. L'EMA et le \code{no\_grad}
sur le GMM brisent la boucle de rétroaction directe.}

\faq{Pourquoi le GMM ne reçoit-il aucun gradient ?}{%
Le GMM est une \emph{cible}, pas un paramètre à optimiser par descente de
gradient. Ses paramètres bougent par mise à jour EMA des statistiques
suffisantes (moyennes/variances pondérées). Laisser un gradient le
traverser créerait un raccourci trivial vers l'effondrement.}

\faq{C'est quoi le \emph{rang effectif} et pourquoi ça aide à choisir une
couche ?}{%
C'est une mesure de « richesse » d'une représentation : si les directions
principales (valeurs singulières) sont nombreuses et équilibrées, le rang
effectif est élevé (l'information est étalée sur beaucoup de dimensions).
Une couche à haut rang effectif tend à mieux séparer les sons, d'où son
choix comme entrée du GMM --- sans avoir besoin d'étiquettes.}

\faq{Pourquoi $K=100$ puis $K=500$ ?}{%
Pour calquer les cardinalités des itérations 1 et 2 de HuBERT. En phase 1,
les features MFCC sont grossières (100 clusters suffisent) ; en phase 2,
les features de l'encodeur sont plus fines et supportent un vocabulaire
plus large (500 clusters), proche du sous-phonémique.}

\subsection{Entraînement pratique}

\faq{Combien de données et de calcul faut-il ?}{%
Le papier pré-entraîne sur $\sim$83\,000 heures (LibriLight + sous-ensemble
anglais de Granary), sur 8 GPU, batch global 192, séquences jusqu'à 15\,s.
C'est un entraînement « gros corpus ». Pour expérimenter, on peut commencer
sur LibriSpeech (960\,h) à échelle réduite.}

\faq{Pourquoi deux \emph{lancements} séparés et pas une seule commande ?}{%
Par simplicité d'ingénierie : la phase 2 a besoin du checkpoint de la phase
1 (\code{--pt\_ckpt}) et active des options différentes (\code{--online\_gmm},
$K$, LR). Conceptuellement c'est \emph{une} trajectoire continue ; en
pratique c'est deux exécutions enchaînées.}

\faq{Que contient le « checkpoint portable » et pourquoi l'appeler ainsi ?}{%
Il bundle l'encodeur online, l'encodeur EMA, les paramètres du GMM et le
numéro de step, dans un simple \file{.pt} indépendant de DeepSpeed. «
Portable » = on peut le recharger sans l'infrastructure d'entraînement,
pour reprendre ou pour extraire des features. C'est lui qu'on chargera côté
Rust pour l'inférence (en ne gardant que l'encodeur).}

\faq{Pourquoi désactiver l'augmentation plus tard en phase 2 ?}{%
Parce que le GMM calcule sa cible sur le signal \emph{propre} alors que
l'encodeur voit le signal \emph{augmenté}. Cela impose une contrainte
d'invariance supplémentaire qui, à ce stade, rend la prédiction masquée
plus difficile que nécessaire. Couper l'augmentation aligne ce que voient
les deux côtés.}

\faq{Pourquoi la loss sur positions \emph{visibles} est-elle calculée mais
pas optimisée ?}{%
Prédire un cluster là où l'encodeur voit déjà l'entrée est trop facile et
apporte des gradients largement redondants. On la garde pour le
\emph{monitoring} (logging) mais on n'optimise que les positions masquées,
comme data2vec / wav2vec~2.0 / HuBERT.}

\subsection{Après l'entraînement : usage}

\faq{Une fois entraîné, qu'est-ce qu'on garde ?}{%
\textbf{Uniquement l'encodeur $f_\phi$}. Le prédicteur, la tête de cluster
et le GMM sont jetés. C'est l'encodeur (51.8M params) qui produit les
représentations utilisées en aval.}

\faq{Comment utilise-t-on S-JEPA pour une tâche concrète (ASR, émotion) ?}{%
On \emph{gèle} l'encodeur et on entraîne une petite tête spécifique à la
tâche sur ses features (protocole SUPERB). L'encodeur n'est pas
fine-tuné ; seule la tête apprend. C'est rapide et révèle la qualité des
représentations.}

\faq{Pourquoi l'inférence en Rust et pas tout en Python ?}{%
Pour la production : Rust offre des latences prévisibles, pas de GIL, une
empreinte mémoire maîtrisée et un déploiement en binaire autonome.
L'entraînement reste en Python (écosystème PyTorch/recherche), l'inférence
critique passe en Rust. On vérifie la parité numérique entre les deux.}

\faq{Comment garantir que le Rust donne les mêmes résultats que le Python ?}{%
Par des \emph{tests de parité} : on sauvegarde des entrées/sorties de
référence depuis Python (frontend CNN, chaque couche transformer, sortie
finale) et on vérifie que l'implémentation Rust reproduit ces tenseurs à
une tolérance numérique près. Points sensibles : LayerNorm en fp32, ordre
des opérations d'attention, weight-norm de l'encodage positionnel.}

\faq{Le GMM est-il nécessaire à l'inférence ?}{%
Non. Le GMM ne sert qu'à fabriquer les \emph{cibles} pendant le
pré-entraînement. À l'inférence on ne fait tourner que l'encodeur ; il n'y
a ni GMM, ni prédicteur, ni masquage.}

\faq{Et le masquage, à l'inférence ?}{%
Désactivé. On donne le signal complet à l'encodeur et on lit ses features
frame par frame. Le masquage est un artefact d'entraînement.}

\subsection{Questions « pièges »}

\faq{Si on jette le GMM et le prédicteur, à quoi servaient-ils ?}{%
Ils ne servent qu'à \emph{façonner} l'encodeur pendant l'entraînement. Le
GMM fournit des cibles douces, le prédicteur force l'encodeur à produire un
contexte assez riche pour qu'on puisse deviner les frames masquées. Une
fois l'encodeur bon, les échafaudages tombent.}

\faq{S-JEPA est-il « juste HuBERT avec un GMM » ?}{%
En esprit oui (« HuBERT à cibles douces dans une architecture JEPA »), mais
trois différences techniques comptent : (1) cibles \emph{douces} GMM au lieu
de dures $k$-means, (2) GMM \emph{en ligne} au lieu de re-clustering
hors-ligne, (3) couche source du GMM choisie \emph{automatiquement} par
rang effectif. C'est la \emph{recette} combinée qui est la contribution,
pas un nouveau bloc d'architecture.}

\faq{Pourquoi S-JEPA bat-il des modèles plus gros sur certaines tâches ?}{%
L'hypothèse du papier : les cibles douces transportent de l'information que
les cibles dures jettent (la masse inter-clusters aux frontières). Sur
36.6\,\% des frames, cette information existe. Le papier ne \emph{prouve}
pas la causalité, mais corrèle cette structure d'incertitude avec
l'efficacité paramétrique.}

\faq{La « switched EMA » est-elle indispensable ?}{%
Le papier la présente comme importante mais reconnaît (Limitations) qu'elle
manque d'une comparaison contrôlée. Le code de référence ne l'implémente
même pas (EMA fixe). Pour la prod, on l'implémentera proprement \emph{et}
on gardera l'option EMA fixe, afin de pouvoir comparer.}

\faq{Tout ceci marche-t-il sur d'autres langues / de l'audio non-vocal ?}{%
Non testé : le papier est entièrement en anglais et le signale comme
limitation. Langues tonales, faibles ressources et audio non-vocal sont
des extensions ouvertes.}

% ============================================================
\section*{Références principales}
% ============================================================
\addcontentsline{toc}{section}{Références principales}
\begin{itemize}[leftmargin=*]
  \item Papier S-JEPA : \file{papers/sources/arXiv-2606.19398v1/main.tex}.
  \item Code de référence :
  \file{papers/sources/s-jepa-main/\{model,train,fit\_gmm\}.py}.
  \item Antécédents : HuBERT (cibles dures, deux itérations),
  wav2vec~2.0 (masquage, encodage positionnel convolutionnel),
  data2vec / BYOL (encodeur EMA), I-JEPA (pattern encodeur--prédicteur).
\end{itemize}

\end{document}
